\documentclass[12pt]{article}

\usepackage[utf8]{inputenc}
\usepackage[T1]{fontenc}
\usepackage[spanish]{babel}
\usepackage{amsmath, amssymb}
\usepackage{geometry}
\usepackage{booktabs}
\usepackage{multirow}
\usepackage{listings}

\lstnewenvironment{Algoritmo}[1]{
    \lstset{
        language=C,
        frame=single,
        basicstyle=\small\ttfamily,
        caption={#1},
        captionpos=t,
        breaklines=true,
        showstringspaces=false,
        literate=
            {á}{{\'a}}1 {é}{{\'e}}1 {í}{{\'\i}}1 {ó}{{\'o}}1 {ú}{{\'u}}1
            {Á}{{\'A}}1 {É}{{\'E}}1 {Í}{{\'I}}1 {Ó}{{\'O}}1 {Ú}{{\'U}}1
            {ñ}{{\~n}}1 {Ñ}{{\~N}}1
    }
}{}

\geometry{margin=2.5cm}

\begin{document}

\section{Sistema de Inventarios (sin espera del cliente)}

\subsection{Demanda diaria}

\begin{description}
    \item[Paso 1: Encontrar la función $f(x)$]

    \begin{center}
    \begin{tabular}{cc}
    \toprule
    $x$ & $f(x)$ \\
    \midrule
    0 & 0.04 \\
    1 & 0.06 \\
    2 & 0.10 \\
    3 & 0.20 \\
    4 & 0.30 \\
    5 & 0.18 \\
    6 & 0.08 \\
    7 & 0.03 \\
    8 & 0.01 \\
    \bottomrule
    \end{tabular}
    \end{center}

    \textbf{Aclaración: esta tabla no es una distribución Poisson.} Para que una variable siga una distribución de Poisson, la media ($\lambda$) y la varianza ($\sigma^2$) tienen que ser exactamente iguales.

    \textbf{Cálculo de la media,} $\text{media} = \sum x \cdot f(x)$, sumando cada valor de $x$ multiplicado por su probabilidad:
    \begin{align*}
    \text{media} &= 0(0.04) + 1(0.06) + 2(0.10) + 3(0.20) + 4(0.30) + 5(0.18) + 6(0.08) + 7(0.03) + 8(0.01) \\
    &= 0 + 0.06 + 0.20 + 0.60 + 1.20 + 0.90 + 0.48 + 0.21 + 0.08 \\
    &= 3.73
    \end{align*}

    \textbf{Cálculo de la varianza,} $\text{varianza} = \sum (x-\text{media})^2 \cdot f(x)$, restando la media a cada $x$, elevando al cuadrado, y multiplicando por su probabilidad:
    \begin{center}
    \small
    \begin{tabular}{ccccc}
    \toprule
    $x$ & $x-3.73$ & $(x-3.73)^2$ & $f(x)$ & $(x-3.73)^2 \cdot f(x)$ \\
    \midrule
    0 & $-3.73$ & 13.9129 & 0.04 & 0.556516 \\
    1 & $-2.73$ & 7.4529  & 0.06 & 0.447174 \\
    2 & $-1.73$ & 2.9929  & 0.10 & 0.299290 \\
    3 & $-0.73$ & 0.5329  & 0.20 & 0.106580 \\
    4 & $0.27$  & 0.0729  & 0.30 & 0.021870 \\
    5 & $1.27$  & 1.6129  & 0.18 & 0.290322 \\
    6 & $2.27$  & 5.1529  & 0.08 & 0.412232 \\
    7 & $3.27$  & 10.6929 & 0.03 & 0.320787 \\
    8 & $4.27$  & 18.2329 & 0.01 & 0.182329 \\
    \bottomrule
    \end{tabular}
    \end{center}
    Sumando la última columna: $0.556516+0.447174+0.299290+0.106580+0.021870+0.290322+0.412232+0.320787+0.182329 = 2.6371$.

    Como la varianza (2.6371) es notablemente menor que la media (3.73) — lo que en estadística se llama subdispersión —, la fórmula de Poisson $\frac{e^{-\lambda}\lambda^x}{x!}$ no se ajusta a este problema. Si se forzara $\lambda=3.73$ en esa fórmula, los valores individuales no coinciden con la tabla: para $x=0$, la tabla da $f(0)=0.04$ pero la fórmula de Poisson daría $e^{-3.73}\approx 0.024$. Por eso esta es una distribución discreta general (tabla de frecuencias empírica), y solo se puede resolver con las sumatorias paso a paso de abajo — no con la fórmula de Poisson.

    \item[Paso 2: Determinar la función acumulada $F(x)$]

    $F(x)$ en cada valor de $x$ es simplemente ir sumando el $f(x)$ de ese valor con el $F(x)$ que ya llevabas acumulado del valor anterior:

    \begin{align*}
    F(0) &= f(0) = 0.04 \\
    F(1) &= F(0) + f(1) = 0.04 + 0.06 = 0.10 \\
    F(2) &= F(1) + f(2) = 0.10 + 0.10 = 0.20 \\
    F(3) &= F(2) + f(3) = 0.20 + 0.20 = 0.40 \\
    F(4) &= F(3) + f(4) = 0.40 + 0.30 = 0.70 \\
    F(5) &= F(4) + f(5) = 0.70 + 0.18 = 0.88 \\
    F(6) &= F(5) + f(6) = 0.88 + 0.08 = 0.96 \\
    F(7) &= F(6) + f(7) = 0.96 + 0.03 = 0.99 \\
    F(8) &= F(7) + f(8) = 0.99 + 0.01 = 1.00
    \end{align*}

    Por eso $F(8)=1.00$: al llegar al último valor ya se sumaron todas las probabilidades de la tabla.

    \begin{center}
    \begin{tabular}{cc}
    \toprule
    $x$ & $F(x)$ \\
    \midrule
    0 & 0.04 \\
    1 & 0.10 \\
    2 & 0.20 \\
    3 & 0.40 \\
    4 & 0.70 \\
    5 & 0.88 \\
    6 & 0.96 \\
    7 & 0.99 \\
    8 & 1.00 \\
    \bottomrule
    \end{tabular}
    \end{center}

    \item[Paso 3: Igualar la función acumulada $F(x) = R$]

    $$ F(x)=R \qquad ; \qquad 0 < R < 1 $$

    Acá ``igualar'' no es despejar una fórmula (como en el caso continuo), sino comparar un $R$ generado contra la tabla de $F(x)$ del Paso 2, hasta encontrar en qué fila deja de cumplirse $R \geq F(x)$. Ejemplo con $R=0.53$:

    \begin{center}
    \begin{tabular}{cccl}
    \toprule
    $x$ & $F(x)$ & ¿$R=0.53 < F(x)$? & \\
    \midrule
    0 & 0.04 & no ($0.53 \geq 0.04$) & sigue \\
    1 & 0.10 & no ($0.53 \geq 0.10$) & sigue \\
    2 & 0.20 & no ($0.53 \geq 0.20$) & sigue \\
    3 & 0.40 & no ($0.53 \geq 0.40$) & sigue \\
    4 & 0.70 & \textbf{sí} ($0.53 < 0.70$) & \textbf{se detiene, $x=4$} \\
    \bottomrule
    \end{tabular}
    \end{center}

    $R=0.53$ es mayor o igual que $F(3)=0.40$ pero menor que $F(4)=0.70$ — por eso $R$ ``iguala'' (queda emparejado) con $x=4$, y no con otro valor.

    \textbf{Demostración de que el tramo asignado es correcto.} Con $0 < R < 1$ generado con \text{gcm}(), por definición de variable uniforme:
    \[
    P(a \leq R < b) = b-a
    \]
    A cada $x_i$ se le asigna el tramo $\big[F(x_{i-1}),\,F(x_i)\big)$. Su ancho es:
    \[
    F(x_i)-F(x_{i-1}) = \sum_{k \leq x_i} f(x_k) - \sum_{k \leq x_{i-1}} f(x_k) = f(x_i)
    \]
    Sustituyendo $a=F(x_{i-1})$, $b=F(x_i)$ en la primera ecuación:
    \[
    P\big(\text{se genera } x_i\big) = P\big(F(x_{i-1}) \leq R < F(x_i)\big) = F(x_i)-F(x_{i-1}) = f(x_i)
    \]
    Con $x_i=4$: $F(3)=0.40$, $F(4)=0.70$, ancho $=0.70-0.40=0.30=f(4)$. Queda demostrado que la probabilidad de generar $x=4$ es exactamente $0.30$, la misma que trae la tabla original — no una aproximación.

    \item[Paso 4: Aplicar la función inversa $x=F^{-1}(R)$]

    Repetir esa comparación fila por fila cada vez que se genera un $R$ sería lento; el Paso 4 arma de una vez la tabla con los tramos ya resueltos (demostrados en el Paso 3), para no tener que comparar contra $F(x)$ cada vez:

    \begin{center}
    \begin{tabular}{ll}
    Si $0.00 \leq R < 0.04$, entonces $x=0$ & Si $0.40 \leq R < 0.70$, entonces $x=4$ \\
    Si $0.04 \leq R < 0.10$, entonces $x=1$ & Si $0.70 \leq R < 0.88$, entonces $x=5$ \\
    Si $0.10 \leq R < 0.20$, entonces $x=2$ & Si $0.88 \leq R < 0.96$, entonces $x=6$ \\
    Si $0.20 \leq R < 0.40$, entonces $x=3$ & Si $0.96 \leq R < 0.99$, entonces $x=7$ \\
                                            & Si $0.99 \leq R \leq 1.00$, entonces $x=8$ \\
    \end{tabular}
    \end{center}

    \item[Paso 5: Propuesta algorítmica para la variable aleatoria]

    En vez de escribir un \texttt{si/sino si} por cada valor, se guardan los valores y su acumulada en dos arreglos y se recorren con un bucle — la misma función sirve para cualquier tabla, sin importar su tamaño:

    \begin{Algoritmo}{}
valores   <- [0, 1, 2, 3, 4, 5, 6, 7, 8]
acumulada <- [0.04, 0.10, 0.20, 0.40, 0.70, 0.88, 0.96, 0.99, 1.00]

entero funcion demandaDiaria(){
    real R <- gcm();
    para i <- 0 hasta longitud(acumulada)-1 hacer
        si (R < acumulada[i]) entonces
            retornar valores[i];
        fin si
    fin para
}
    \end{Algoritmo}
\end{description}

\subsection{Tiempo de entrega}

\begin{description}
    \item[Paso 1: Encontrar la función $f(x)$]
    

    \begin{center}
    \begin{tabular}{cc}
    \toprule
    $x$ (días) & $f(x)$ \\
    \midrule
    1 & 0.25 \\
    2 & 0.50 \\
    3 & 0.20 \\
    4 & 0.05 \\
    \bottomrule
    \end{tabular}
    \end{center}

    \item[Paso 2: Determinar la función acumulada $F(x)$]

    \begin{align*}
    F(1) &= f(1) = 0.25 \\
    F(2) &= F(1) + f(2) = 0.25 + 0.50 = 0.75 \\
    F(3) &= F(2) + f(3) = 0.75 + 0.20 = 0.95 \\
    F(4) &= F(3) + f(4) = 0.95 + 0.05 = 1.00
    \end{align*}

    \begin{center}
    \begin{tabular}{cc}
    \toprule
    $x$ (días) & $F(x)$ \\
    \midrule
    1 & 0.25 \\
    2 & 0.75 \\
    3 & 0.95 \\
    4 & 1.00 \\
    \bottomrule
    \end{tabular}
    \end{center}

    \item[Paso 3: Igualar la función acumulada $F(x) = R$]

    $$ F(x)=R \qquad ; \qquad 0 < R < 1 $$

    Igual que con la demanda: se compara el $R$ generado contra la tabla de $F(x)$ hasta encontrar la primera fila donde $R < F(x)$. Ejemplo con $R=0.42$: $R \geq F(1)=0.25$ pero $R < F(2)=0.75$, entonces el tiempo de entrega es $x=2$ días.

    \item[Paso 4: Aplicar la función inversa $x=F^{-1}(R)$]

    Con los tramos ya resueltos, sin tener que comparar contra $F(x)$ cada vez:

    \begin{center}
    \begin{tabular}{l}
    Si $0.00 \leq R < 0.25$, entonces $x=1$ día \\
    Si $0.25 \leq R < 0.75$, entonces $x=2$ días \\
    Si $0.75 \leq R < 0.95$, entonces $x=3$ días \\
    Si $0.95 \leq R \leq 1.00$, entonces $x=4$ días \\
    \end{tabular}
    \end{center}

    \item[Paso 5: Propuesta algorítmica para la variable aleatoria]

    Misma idea: valores y acumulada en arreglos, un bucle genérico:

    \begin{Algoritmo}{}
valores   <- [1, 2, 3, 4]
acumulada <- [0.25, 0.75, 0.95, 1.00]

entero funcion tiempoEntrega(){
    real R <- gcm();
    para i <- 0 hasta longitud(acumulada)-1 hacer
        si (R < acumulada[i]) entonces
            retornar valores[i];
        fin si
    fin para
}
    \end{Algoritmo}
\end{description}

\subsection{Algoritmo de simulación del inventario}

\begin{Algoritmo}{}
real funcion simularInventario(entero q, entero R, entero dias){
    entero inventario, faltanteAcum, ordenes, diaLlegada, hayPedido;
    real sumaProm, costoOrdenar, costoInventario, costoFaltante;

    inventario <- 15;
    faltanteAcum <- 0;
    ordenes <- 0;
    sumaProm <- 0;
    hayPedido <- falso;

    para dia <- 1 hasta dias hacer
        si (hayPedido) && (dia == diaLlegada) entonces
            inventario <- inventario + q;
            hayPedido <- falso;
        fin si

        invInicial <- inventario;
        demanda <- demandaDiaria();

        si (demanda <= inventario) entonces
            inventario <- inventario - demanda;
            promedioDia <- (invInicial + inventario) / 2;
        sino
            faltanteAcum <- faltanteAcum + (demanda - inventario);
            // el inventario se agota antes de terminar el dia: el promedio no es
            // (invInicial+0)/2, es el promedio de la parte con stock (invInicial/2)
            // multiplicado por la fraccion del dia que duro esa parte (invInicial/demanda)
            si (invInicial > 0) entonces
                promedioDia <- (invInicial / 2) * (invInicial / demanda);
            sino
                promedioDia <- 0;
            fin si
            inventario <- 0;
        fin si

        sumaProm <- sumaProm + promedioDia;

        si (inventario <= R) && (!hayPedido) entonces
            diaLlegada <- dia + tiempoEntrega();
            hayPedido <- verdadero;
            ordenes <- ordenes + 1;
        fin si
    fin para

    costoOrdenar <- ordenes * 50;
    costoInventario <- sumaProm * (26.0 / dias);
    costoFaltante <- faltanteAcum * 25;

    retornar costoOrdenar + costoInventario + costoFaltante;
}
\end{Algoritmo}

\subsection{Simulación manual ($q=30$, $R=10$)}

\textbf{Qué significa ``colocar una orden'' y cómo se usa el tiempo de entrega:}

Cuando el inventario al final del día queda en $R=10$ o menos, y no hay ya un pedido en camino, se coloca una orden de $q=30$ unidades. La cantidad que llega es \textbf{siempre $q=30$}, sin importar cuánta era la demanda ese día ni en cuánto quedó el inventario — es la política de ``lote constante'': lo único que varía entre una orden y otra es cuándo llega (el tiempo de entrega), nunca cuánto llega. Colocar la orden son dos cosas concretas:
\begin{enumerate}
    \item Se genera \textbf{cuándo va a llegar}: se saca un $R$ nuevo (distinto del que generó la demanda de ese día) y se lo pasa por la tabla de tiempo de entrega (Paso 4 de la sección anterior).
    \item Se calcula el día de llegada $=$ día en que se colocó la orden $+$ tiempo de entrega obtenido.
\end{enumerate}
Mientras se espera esa orden, no se coloca ninguna otra — no tiene sentido pedir de nuevo antes de que llegue lo que ya se pidió.

Ejemplo con la tabla de abajo: el día 2 el inventario terminó en 6, que es $\leq R=10$, y no había pedido pendiente $\Rightarrow$ se coloca la orden 1 (columna ``Orden'' marca ``1'' en esa fila). Para saber cuándo llega, se genera $R=0.42$ para el tiempo de entrega: cae en el tramo $0.25 \leq R < 0.75 \Rightarrow 2$ días. Día de llegada $= 2+2 = 4$.

Los días 2 y 3 el inventario sigue bajando con la demanda normal — todavía no llegó nada. Recién el día 4, \textbf{antes} de generar la demanda de ese día, se suma $q=30$ al inventario que había quedado el día 3 (4 unidades): $4+30=34$. Por eso en la tabla el ``Inventario inicial'' del día 4 es 34 y no 4. Después de sumar el pedido, el día 4 sigue como cualquier otro día: se genera su propia demanda y se resta normalmente.

La columna ``Orden'' de la tabla marca únicamente el día en que \emph{se coloca} la orden (día 2 $\to$ ``1'', día 10 $\to$ ``2''); no marca el día en que llega — eso queda anotado aparte, debajo de la tabla.

\textbf{Inventario promedio en días con faltante (día 11):} el inventario empieza el día en 4, pero la demanda es 8 — no alcanza, y el inventario llega a 0 \emph{antes} de que termine el día (se agota apenas se cubre una fracción de la demanda). Por eso el promedio de ese día no es $(4+0)/2=2.0$; usando la misma fórmula del Ejemplo 5.5 de \texttt{cap5\_simulacion.tex} (Tabla 5.7):
\[
\text{promedio día 11} = \frac{4}{2} \times \frac{4}{8} = 2 \times 0.5 = 1.0
\]
El día 12 el inventario ya empieza en 0 (invInicial$=0$), así que el promedio de ese día es directamente $0$ — no hay nada que promediar.

Inventario inicial: 15 unidades.

\begin{center}
\small
\begin{tabular}{cccccccc}
\toprule
\multirow{2}{*}{Día} & Inventario & Número & \multirow{2}{*}{Demanda} & Inventario & \multirow{2}{*}{Faltante} & \multirow{2}{*}{Orden} & Inventario \\
 & inicial & aleatorio & & final & & & diario promedio \\
\midrule
1  & 15 & 0.53 & 4 & 11 &   &   & 13.0 \\
2  & 11 & 0.81 & 5 & 6  &   & 1 & 8.5  \\
3  & 6  & 0.15 & 2 & 4  &   &   & 5.0  \\
4  & 34 & 0.68 & 4 & 30 &   &   & 32.0 \\
5  & 30 & 0.09 & 1 & 29 &   &   & 29.5 \\
6  & 29 & 0.77 & 5 & 24 &   &   & 26.5 \\
7  & 24 & 0.95 & 6 & 18 &   &   & 21.0 \\
8  & 18 & 0.31 & 3 & 15 &   &   & 16.5 \\
9  & 15 & 0.61 & 4 & 11 &   &   & 13.0 \\
10 & 11 & 0.98 & 7 & 4  &   & 2 & 7.5  \\
11 & 4  & 0.99 & 8 & 0  & 4 &   & 1.0  \\
12 & 0  & 0.45 & 4 & 0  & 4 &   & 0.0  \\
13 & 30 & 0.22 & 3 & 27 &   &   & 28.5 \\
\bottomrule
\end{tabular}
\end{center}

Orden 1 (día 2): $R=0.42 \Rightarrow$ tiempo de entrega $=2$ días $\Rightarrow$ llega día 4.
Orden 2 (día 10): $R=0.88 \Rightarrow$ tiempo de entrega $=3$ días $\Rightarrow$ llega día 13.

\textbf{De dónde sale cada número del costo:} los tres costos se leen directo de las columnas de la tabla de arriba, no de ningún dato nuevo.

\begin{itemize}
    \item \textbf{2} (para el costo de ordenar) = cuántas veces aparece algo en la columna ``Orden'' = 1 vez en el día 2 + 1 vez en el día 10 = 2 órdenes en total.
    \item \textbf{202.0} (para el costo de inventario) = la suma de \emph{toda} la columna ``Inventario diario promedio'', los 13 días (el día 11 usa la fórmula corregida, $1.0$, no $2.0$):
    \[
    13.0+8.5+5.0+32.0+29.5+26.5+21.0+16.5+13.0+7.5+1.0+0.0+28.5 = 202.0
    \]
    \item \textbf{8} (para el costo de faltante) = la suma de \emph{toda} la columna ``Faltante'': $4$ (día 11) $+$ $4$ (día 12) $= 8$; en los demás días esa columna está vacía porque no hubo faltante.
\end{itemize}

Con esos tres números se aplican directamente las fórmulas de costo del \texttt{Algoritmo de simulación del inventario} (\texttt{costoOrdenar <- ordenes * 50}, etc.):

\begin{center}
\begin{tabular}{cccc}
\toprule
Costo de ordenar & Costo de inventario & Costo de faltante & Costo total \\
\midrule
$2(50)=\$100$ & $202.0(26/260)=\$20.20$ & $8(25)=\$200$ & \$320.20 \\
\bottomrule
\end{tabular}
\end{center}

El costo total (\$320.20) es simplemente la suma de esos tres: $100+20.20+200=320.20$.

\section{Preguntas pendientes para el docente (p3 y p4)}

\begin{enumerate}
    \item ¿Se revisa el inventario todos los días para decidir si hay que pedir más mercadería, o se revisa solo cada cierto tiempo (por ejemplo, cada 3 días)? Nosotros asumimos que se revisa todos los días.

    \item ¿Cómo se debe contar cuánto inventario ``tuviste guardado'' en el día? Nosotros usamos el promedio entre lo que había al empezar el día y lo que quedó al final.

    \item Probamos varios pares $(q,R)$ y varios dieron casi el mismo costo — no hay un único par ``perfecto''. ¿El docente espera un $q$ y $R$ exactos, o alcanza con reportar el rango de valores que dan el mejor resultado?

    \item Para buscar el mejor $(q,R)$ probamos varios pares y nos quedamos con el de menor costo. ¿Alcanza con eso, o se espera un método de búsqueda más formal?

    \item (Solo p4) Si dos clientes están esperando su pedido y llega mercadería que no alcanza para los dos, ¿a quién se le entrega primero? Nosotros decidimos que al que lleva más tiempo esperando.

    \item (Solo p4) ¿Cómo se debe contar si un cliente ``se cansó de esperar''? Nosotros lo revisamos día por día: si el pedido no llega antes de que se le acabe la paciencia al cliente, esa venta se pierde.

    \item Para p3 y p4 solo hicimos la implementación en Java (sin Excel). ¿Está bien así, o hace falta Excel también para estos dos problemas?

    \item ¿Hay una forma o símbolos específicos que el docente pide para el diagrama de flujo de la simulación de inventarios?

    \item ¿Cuántos días o años hace falta simular para confiar en el resultado? Nosotros simulamos hasta 1000 años para que el promedio ya no cambie.
\end{enumerate}

\end{document}
